\scanpage{7}\mspage{5}

façon équivalente, il existe un $v\in \unsure{E|\mathcal E}^{2n-2i}$, tel que le morphisme induit $v:M^i- M^{2n-i}$ soit un isomorphisme.

On peut transformer ceci de la façon suivante :
\[
 \left\{\begin{aligned}
 &\exists\ v_i\in E^{2n-2i},\ w_i\in E^{-(2n-2i)},\quad\text{tels que l'on ait}\\
 &(w_i v_i-1)\pi_i=0,\qquad(v_i w_i-1)\pi_{2n-i}=0
 \end{aligned}\right.
 \tag*{$(4.4.)_i$}
\]
Or, sans plus supposer que $E^*$ se réalise à l'aide d'un $M^*$, on voit aussitôt que $(4.4.)_i$ implique $(3.1)_i$ et $(3.2)_i$. En effet, quitte à remplacer $v,w$ par $v\pi_i=\pi_{2n-i}v$, $w\pi_{2n-i}=\pi_i w$, on \unsure{va supposer|les suppose|peut supposer} $v,w$ dans $E_i^{2n-i}$, $E_{2n-i}^i$. Si alors \deleted{\unsure{$Z$}} $u\in E^{(0)}$ est tel que $E_i^{2n-i}u=0$, donc $vu=0$ on aura $wvu\pi_i\deleted{\unsure{0}}=0$ i.e. $wv\pi_i u=0$ donc $\pi_i u=0$ (car $wv\pi_i=\pi_i$). De $=$ pour $(3.2)_i$, en utilisant la deuxième relation $(\unsure{4.4}.)_i$.

\margin{\uline{cf 5. v. p \uline{5bis}}}

\deleted{\unsure{5}}6. Nous allons étudier de façon plus détaillée un cas où les conditions (4.4.) sont satisfaites. Pour ceci, partons encore du cas particulier $M^*$, \inserted{$E^*$ sous-algèbre graduée de $\operatorname{End}(M^*)$,} et supposons qu'on ait un
\[
 L\in E^{(2)}
 \tag*{(6.0)}
\]
qui ait la propriété que
% archive sheet number: 6
